\section{Stock Forecast Evaluation}\label{sec:stock_validation_supp}

The stock evaluation separated a check of the existing generator from the comparison of new benchmarks. The saved GS and LLY marginals used daily steps of $1/252$ year and annualized excess log returns. Reconstructing the first-origin forecasts reproduced the saved stock means at both one and five sessions. The original simulator drew the first return state from the stationary distribution, without conditioning on recent returns. In the initialization check, we inferred the current state from available daily returns and advanced it before generating the first future return. The inference tracked both the return state and the number of forced tail-state emissions remaining in a jump episode. Poisson durations retained all but at most $10^{-12}$ probability, with the remaining mass assigned to the final duration. Missing daily observations advanced the latent distribution without an observation likelihood; an accumulated multi-day return was not treated as a daily observation. The fixed pilot growth shift was removed from observed annualized returns when calculating likelihoods, consistent with the simulated observation model. Table~\ref{tab:stock_initialization} compares this conditional initialization with the original start on the same 2026 endpoints. Both used three 10{,}000-path replicates, with seeds separate from the subsequent four-method benchmark. The original return-state distributions mixed rapidly under the fitted transitions, so the initial-state adjustment had little effect on five-session central forecasts. The broader exploratory checks, including a fixed 20-return volatility adjustment and Gaussian random walks, are retained in the saved experiment outputs.

\begin{table}[H]
\centering
\caption{Stock forecast errors under original and history-conditioned initialization on identical 2026 endpoints. Each ticker has 21 one-session and 17 five-session origins. MAE scores the forecast median; errors are dollars per share and average three 10{,}000-path replicates. Filtered denotes initialization from available past returns, including inference of remaining jump duration.}\label{tab:stock_initialization}
\resizebox{\textwidth}{!}{\begin{tabular}{lrrrrr}
\toprule
Stock & Horizon & Original MAE & Filtered MAE & Original CRPS & Filtered CRPS \\
\midrule
GS & 1 & 10.64 & 10.76 & 7.99 & 8.41 \\
GS & 5 & 16.28 & 16.35 & 12.99 & 13.19 \\
LLY & 1 & 19.10 & 19.09 & 14.73 & 14.64 \\
LLY & 5 & 49.58 & 49.24 & 35.69 & 35.28 \\
\bottomrule
\end{tabular}
}
\end{table}

The benchmark comparison retained the original generator and tested one adaptive variance rule and one directional model against it. For each annual validation block from 2019 through 2024, regression training responses ended before the validation year and validation endpoints remained within it. The EWMA variance began with the sample variance of the first sixty training daily returns and then updated using only returns observed by each origin. Decay selection minimized $\tfrac12[\log(h a_t)+y_{t,h}^2/(h a_t)]$, where $y_{t,h}$ was the observed cumulative log return over $h$ sessions. We first averaged this loss within each ticker, validation year, and horizon, then gave the resulting groups equal weight. With the selected decay, the regression fitted the next daily return divided by origin volatility. Its four features were the stock's last daily return, its last five- and twenty-return sums, and SPY's last five-return sum. Return sums were divided by the corresponding origin volatility and square root of window length; feature columns were further divided by their training root mean square. The ridge objective was mean squared normalized one-session error plus the penalty times the squared coefficient norm. There was no intercept. Penalty selection used $(y_{t,h}-h\widehat\mu_t)^2/(h a_t)$, with the same group weighting as decay selection. Table~\ref{tab:stock_selection} gives every candidate's validation loss. The selected zero-drift limit set all four coefficients to zero for both tickers. Final coefficients and feature scales were fitted using 2014--2024 only and saved before evaluating test forecasts. The two selection losses answer different questions and are not comparable in magnitude.

\begin{table}[H]
\centering
\caption{Earlier-data selection of adaptive-volatility and directional-model settings. Losses average ticker, year, and horizon groups equally across 2019--2024 validation blocks. Lower values are better within each selection task. EWMA decay $0.97$ and the zero-drift limit were selected; 2025 and 2026 outcomes did not enter these choices.}\label{tab:stock_selection}
\begin{tabular}{llr}
\toprule
Selection & Setting & Validation loss \\
\midrule
EWMA decay & 0.80 & -2.921492 \\
EWMA decay & 0.90 & -3.009741 \\
EWMA decay & 0.94 & -3.044490 \\
EWMA decay & 0.97 & -3.065558 \\
EWMA decay & 0.99 & -3.064379 \\
Ridge penalty & 0.01 & 1.146523 \\
Ridge penalty & 0.1 & 1.143850 \\
Ridge penalty & 1 & 1.138475 \\
Ridge penalty & 10 & 1.136428 \\
Ridge penalty & 100 & 1.136142 \\
Ridge penalty & 1000 & 1.136112 \\
Ridge penalty & Zero drift & 1.136109 \\
\bottomrule
\end{tabular}

\end{table}

Historical close series for GS, LLY, and SPY shared the same recorded session grid, with 2{,}767 closes per ticker in 2014--2024 and 250 in 2025. The 2025 test updated lag features and variance using information through each origin while keeping fitted parameters fixed. The 2026 series began in April after a data gap. Its variance restarted from the 2014--2024 sample estimate and updated through the available April--July daily returns before evaluation. December 2025 and April 2026 were not joined into a daily return. Within 2026, both adjacent exchange-session closes were required to calculate a daily return; missing observations left variance unchanged. Feature windows used the last required number of valid daily returns, which could span more sessions when observations were missing. The oldest feature input at an evaluated 2026 origin was four calendar days old, whereas the 2025 features had no such gap. Every method used the same observed starting price and endpoint, and none used option quotes. Tables~\ref{tab:stock_validation_five} and~\ref{tab:stock_validation_one} retain all four candidates and both evaluation periods. The directional and adaptive rows agree exactly because zero drift was selected. Their small median-error differences from the unchanged-price benchmark arise from Monte Carlo estimation of a distribution whose population median equals the origin price. They are not a learned directional effect. Replicate scores, date-level comparisons, input hashes, and settings are retained with the experiment. Overlapping horizons imply dependent observations, and the three simulation seeds measure simulation variation rather than three independent market samples.

\begin{table}[H]
\centering
\caption{Complete five-session stock forecast scores for the four selected candidates. Each ticker has 245 origins in 2025 and 17 in August--September 2026. MAE scores medians and RMSE scores means. Monetary quantities are dollars per share; coverage and width describe central 90\% intervals. Scores average dates within each of three 10{,}000-path replicates and then average replicates. Directional equals Adaptive because earlier-data selection chose zero drift.}\label{tab:stock_validation_five}
\resizebox{\textwidth}{!}{\begin{tabular}{lllrrrrr}
\toprule
Year & Stock & Model & Median MAE & Mean RMSE & CRPS & Coverage (\%) & Width \\
\midrule
2025 & GS & JumpHMM & 22.31 & 28.42 & 16.22 & 84.4 & 84.71 \\
2025 & GS & Unchanged & 22.61 & 28.75 & 22.61 & -- & -- \\
2025 & GS & Adaptive & 22.62 & 28.57 & 16.32 & 89.8 & 94.86 \\
2025 & GS & Directional & 22.62 & 28.57 & 16.32 & 89.8 & 94.86 \\
2025 & LLY & JumpHMM & 39.10 & 50.16 & 29.60 & 69.9 & 100.40 \\
2025 & LLY & Unchanged & 39.37 & 50.36 & 39.37 & -- & -- \\
2025 & LLY & Adaptive & 39.35 & 50.23 & 28.57 & 84.4 & 151.37 \\
2025 & LLY & Directional & 39.35 & 50.23 & 28.57 & 84.4 & 151.37 \\
2026 & GS & JumpHMM & 16.25 & 21.71 & 13.03 & 100.0 & 129.05 \\
2026 & GS & Unchanged & 15.51 & 20.71 & 15.51 & -- & -- \\
2026 & GS & Adaptive & 15.48 & 21.07 & 15.58 & 100.0 & 176.18 \\
2026 & GS & Directional & 15.48 & 21.07 & 15.58 & 100.0 & 176.18 \\
2026 & LLY & JumpHMM & 49.44 & 58.63 & 35.60 & 70.6 & 147.26 \\
2026 & LLY & Unchanged & 49.37 & 58.17 & 49.37 & -- & -- \\
2026 & LLY & Adaptive & 49.31 & 58.35 & 33.81 & 94.1 & 191.38 \\
2026 & LLY & Directional & 49.31 & 58.35 & 33.81 & 94.1 & 191.38 \\
\bottomrule
\end{tabular}
}
\end{table}

\begin{table}[H]
\centering
\caption{Complete one-session stock forecast scores for the four selected candidates. Each ticker has 249 origins in 2025 and 21 in August--September 2026. Units, point summaries, interval definitions, and aggregation follow Supplementary Table~\ref{tab:stock_validation_five}.}\label{tab:stock_validation_one}
\resizebox{\textwidth}{!}{\begin{tabular}{lllrrrrr}
\toprule
Year & Stock & Model & Median MAE & Mean RMSE & CRPS & Coverage (\%) & Width \\
\midrule
2025 & GS & JumpHMM & 9.10 & 12.26 & 6.66 & 87.8 & 34.67 \\
2025 & GS & Unchanged & 9.13 & 12.28 & 9.13 & -- & -- \\
2025 & GS & Adaptive & 9.13 & 12.27 & 6.74 & 93.4 & 42.44 \\
2025 & GS & Directional & 9.13 & 12.27 & 6.74 & 93.4 & 42.44 \\
2025 & LLY & JumpHMM & 14.03 & 20.42 & 10.67 & 77.5 & 40.23 \\
2025 & LLY & Unchanged & 14.07 & 20.44 & 14.07 & -- & -- \\
2025 & LLY & Adaptive & 14.06 & 20.43 & 10.63 & 93.3 & 67.61 \\
2025 & LLY & Directional & 14.06 & 20.43 & 10.63 & 93.3 & 67.61 \\
2026 & GS & JumpHMM & 10.67 & 14.34 & 8.00 & 90.5 & 52.39 \\
2026 & GS & Unchanged & 10.59 & 14.23 & 10.59 & -- & -- \\
2026 & GS & Adaptive & 10.60 & 14.30 & 8.72 & 100.0 & 77.21 \\
2026 & GS & Directional & 10.60 & 14.30 & 8.72 & 100.0 & 77.21 \\
2026 & LLY & JumpHMM & 19.11 & 25.74 & 14.73 & 74.6 & 58.00 \\
2026 & LLY & Unchanged & 19.01 & 25.79 & 19.01 & -- & -- \\
2026 & LLY & Adaptive & 18.98 & 25.78 & 14.47 & 76.2 & 83.93 \\
2026 & LLY & Directional & 18.98 & 25.78 & 14.47 & 76.2 & 83.93 \\
\bottomrule
\end{tabular}
}
\end{table}
\FloatBarrier
